\section{Experimental Evaluation: Deterministic Forecasting}
\label{sec:deterministic}

We now present a comprehensive evaluation of the model variants introduced in Section~\ref{sec:neural_ode}. Our experiments employ rigorous out-of-sample evaluation, training on historical data and testing on a completely held-out future period. This design mirrors operational forecasting conditions and reveals which architectural choices lead to genuine predictive skill versus overfitting to training-period patterns.

\subsection{Experimental Protocol}

Our evaluation follows a strict temporal separation between training and test data. The training period spans January 1979 through December 2001, comprising 276 monthly observations over 23 years. The test period spans January 2002 through December 2022, comprising 252 monthly observations over 21 years that are held out entirely during model development. This temporal split ensures that test performance reflects true out-of-sample generalization rather than interpolation within the training distribution.

We evaluate forecasts at lead times from 1 to 21 months, covering the operationally-relevant horizon for ENSO prediction. At each initialization time in the evaluation period, we integrate the model forward and compare the predicted Ni\~no3.4 index against observations. Two complementary metrics quantify forecast skill. The root mean squared error (RMSE) measures the typical magnitude of forecast errors in degrees Celsius, with lower values indicating better performance. The anomaly correlation coefficient (ACC) measures the linear association between forecasts and observations, ranging from $-1$ to $+1$ with higher values indicating better performance. We also compute a combined score (normalized ACC minus normalized RMSE) for ranking purposes.

All models are trained using the procedure described in Section~\ref{sec:neural_ode}: Adam~\cite{kingma2015adam} optimization with learning rate $10^{-3}$, weight decay~\cite{krogh1992simple} $10^{-4}$, and early stopping~\cite{prechelt1998early} based on validation loss. We train for up to 1500 epochs, with the validation set comprising forecasts initialized from the final 20\% of the training period.

\subsection{Overall Performance Comparison}

Table~\ref{tab:top10_outsample} presents the performance of the top 10 model variants along with key baselines. The results reveal several striking patterns that inform our understanding of physics-informed machine learning for climate prediction.

\begin{table}[htbp]
\centering
\caption{Model performance ranked by out-of-sample skill. Training period: 1979--2001; test period: 2002--2022. Models are ranked by a combined score of normalized ACC and RMSE.}
\label{tab:top10_outsample}
\begin{tabular}{clcccc}
\toprule
\textbf{Rank} & \textbf{Model} & \textbf{Test ACC} & \textbf{Test RMSE} & \textbf{Train RMSE} & \textbf{Gap} \\
\midrule
1 & NXRO-Res & 0.628 & 0.567°C & 0.474°C & 0.093°C \\
2 & NXRO-Graph & 0.616 & 0.583°C & 0.536°C & 0.047°C \\
3 & NXRO-Attentive & 0.613 & 0.586°C & 0.542°C & 0.044°C \\
4 & NXRO-RO+Diag & 0.611 & 0.618°C & 0.490°C & 0.128°C \\
5 & NXRO-Linear & 0.611 & 0.588°C & 0.543°C & 0.045°C \\
6 & NXRO-NeuralODE & 0.602 & 0.588°C & 0.522°C & 0.066°C \\
7 & XRO (baseline) & 0.596 & 0.605°C & 0.493°C & 0.112°C \\
\midrule
40 & NXRO-Transformer & 0.482 & 0.651°C & 0.678°C & $-$0.027°C \\
\bottomrule
\end{tabular}
\end{table}

The residual neural model (NXRO-Res) achieves the best out-of-sample performance, with a test RMSE of 0.567°C and test ACC of 0.628. This represents a 6.3\% improvement in RMSE and 5.4\% improvement in ACC compared to the XRO baseline. Notably, the residual model achieves this performance with a simpler architecture than XRO, omitting the polynomial nonlinearities entirely and relying instead on the MLP to capture any nonlinear dynamics present in the data.

Figure~\ref{fig:top10_rmse} displays RMSE as a function of forecast lead time for the top models. The residual model maintains its advantage across all lead times, with particularly pronounced improvements at longer leads (12--21 months) where error accumulation is most severe. This suggests that the learned residual captures dynamics relevant for extended forecasting that the polynomial basis misses or misspecifies.

\begin{figure}[htbp]
    \centering
    \includegraphics[width=0.95\textwidth]{top10_rmse_test_out_of_sample.png}
    \caption{Out-of-sample RMSE as a function of forecast lead time. The residual model (NXRO-Res) achieves lowest error at all leads. All hybrid models outperform the XRO baseline, while the pure transformer performs worst.}
    \label{fig:top10_rmse}
\end{figure}

\subsection{Analysis of Alternative Architectures}

The ranking in Table~\ref{tab:top10_outsample} encodes rich information about which architectural choices matter for climate forecasting. We now analyze each class of models to extract generalizable insights.

\subsubsection{Graph-Based Models}

The graph-based variant (NXRO-Graph) ranks second overall, demonstrating that structured sparsity can effectively regularize the nonlinear component. By restricting interactions to follow the adjacency pattern derived from XRO's linear operator, the graph model prevents learning of spurious correlations that might arise from allowing all-to-all coupling. The test RMSE of 0.583°C represents a 3.6\% improvement over XRO.

Interestingly, the graph model exhibits the smallest generalization gap among the top performers (0.047°C between train and test RMSE). This suggests that the graph structure provides strong regularization without sacrificing expressiveness. The sparse connectivity acts as an architectural prior encoding the belief that climate modes interact through specific teleconnection pathways rather than through dense, unstructured coupling.

The success of the graph variant has implications for interpretability. Unlike the MLP residual, the graph convolution produces coupling strengths that can be inspected and compared to physical understanding of teleconnections. Future work could leverage this structure to discover previously unknown interaction pathways.

\subsubsection{Attention-Based Models}

The attention-based variant (NXRO-Attentive) ranks third with test RMSE of 0.586°C. Unlike the graph model, which uses fixed connectivity, the attention mechanism learns state-dependent coupling weights that vary with the current climate state. This allows the model to capture regime-dependent dynamics---for instance, different interaction patterns during El Ni\~no versus La Ni\~na conditions.

The masked attention architecture proves crucial. When we remove the attention mask and allow unrestricted self-attention (equivalent to a transformer applied to variables as tokens), performance degrades substantially. The mask, which restricts peripheral modes from directly influencing core ENSO variables, prevents the attention mechanism from learning spurious dependencies that happen to correlate in the training period but do not generalize.

The attention model's generalization gap (0.044°C) is comparable to the graph model, suggesting that the mask provides similar regularization benefits to the graph's sparse connectivity. Both approaches encode domain knowledge about which interactions are physically plausible, using this prior to constrain the hypothesis space.

\subsubsection{Physics-Structured Models}

The gradient-trained physics model (NXRO-RO+Diag) replicates XRO's polynomial structure but optimizes all coefficients jointly via gradient descent rather than variable-by-variable closed-form fitting. This variant ranks fourth with test RMSE of 0.618°C, slightly worse than XRO's 0.605°C despite having identical functional form.

This result initially seems paradoxical: end-to-end optimization should find at least as good a solution as the closed-form fit, if not better. The explanation lies in the generalization gap. NXRO-RO+Diag achieves substantially better training RMSE (0.490°C vs 0.493°C) but worse test RMSE, exhibiting a gap of 0.128°C---the largest among the top 10 models. The gradient-based optimization finds coefficient configurations that fit the training trajectories more closely but generalize less well to the test period.

This overfitting occurs because the polynomial basis has many degrees of freedom (five RO monomials with seasonal coefficients, plus diagonal terms for each variable), and the multi-step loss provides a rich signal that can be exploited to fit training-period idiosyncrasies. XRO's closed-form fit, which optimizes one-step predictions independently, may act as implicit regularization by preventing the exploitation of multi-step correlations.

\subsubsection{Pure Linear Models}

The pure linear variant (NXRO-Linear) ranks fifth with test RMSE of 0.588°C, outperforming both XRO and the gradient-trained physics model despite having no nonlinear terms at all. This remarkable result underscores the regularization value of simplicity when data is limited.

The linear model's success suggests that much of ENSO predictability arises from the seasonally-modulated linear dynamics captured by the operator $\bL(t)$, with nonlinear effects contributing less to forecast skill than commonly assumed. The polynomial nonlinearities in XRO may primarily fit amplitude-dependent fluctuations in the training period that do not persist into the test period.

Alternatively, the gradient-based optimization of the linear operator may find a different coefficient configuration than XRO's closed-form fit---one that implicitly compensates for the missing nonlinearities by adjusting the linear coupling strengths. Disentangling these effects would require comparing the fitted operators directly, which we leave for future work.

\subsubsection{Pure Neural Models}

The pure neural ODE (NXRO-NeuralODE) ranks sixth with test RMSE of 0.588°C, matching the linear model despite having substantially more capacity. This parity suggests that without the linear operator's inductive bias, the MLP struggles to learn the seasonally-structured dynamics from limited data.

The neural ODE receives seasonal information through the Fourier features $\bphi(t)$, but this provides weaker structure than the explicit Fourier expansion of the linear operator. The MLP must learn from scratch that dynamics vary smoothly with season, rather than having this structure encoded architecturally. With only 276 training months spanning 23 annual cycles, this learning problem is severely underdetermined.

The neural ODE's moderate generalization gap (0.066°C) indicates that it neither overfits dramatically nor achieves the excellent training fit of the physics-structured models. It represents a reasonable but not optimal tradeoff between flexibility and regularization.

\subsubsection{Transformer Models}

The transformer-based variant (NXRO-Transformer) ranks 40th out of 43 models evaluated, with test RMSE of 0.651°C---substantially worse than the XRO baseline's 0.605°C. This represents a 7.6\% degradation in RMSE and a 19\% degradation in ACC. The transformer performs poorly even on training data (train RMSE 0.678°C), indicating fundamental rather than overfitting-related failure.

This negative result carries important implications for practitioners. Transformers~\cite{vaswani2017attention} have achieved remarkable success in natural language processing and computer vision~\cite{dosovitskiy2020image}, leading to enthusiasm for their application across domains. However, these successes typically involve massive datasets (millions to billions of examples) and sequential structure that transformers are specifically designed to exploit.

Climate forecasting presents a fundamentally different challenge. The dataset comprises only 276 training months, and the state vector of 10 variables does not exhibit the sequential structure of text or spatial structure of images. While hybrid architectures combining CNNs and transformers have shown promise for ENSO forecasting~\cite{lyu2024resonet,patil2023deep}, these approaches typically leverage spatiotemporal data from full SST fields rather than derived indices. The self-attention mechanism, which computes all-to-all interactions between variables at each time step, has many parameters relative to the available data and lacks the inductive biases needed to generalize from limited examples.

The contrast between NXRO-Attentive (rank 3, 0.586°C) and NXRO-Transformer (rank 40, 0.651°C) is particularly instructive. Both use attention mechanisms, but NXRO-Attentive retains the physics-informed linear operator and employs masked attention to restrict interactions. The 0.065°C gap between these architectures---larger than the gap between any other adjacent models in the top 10---demonstrates that physics-informed structure provides essential regularization that cannot be recovered through increased model capacity.

\subsection{The Generalization Gap}

Table~\ref{tab:generalization} analyzes the relationship between training performance and generalization for representative models. The generalization gap, defined as the difference between test and training RMSE, quantifies how much performance degrades when moving from in-sample to out-of-sample evaluation.

\begin{table}[htbp]
\centering
\caption{Generalization analysis for selected models. The gap measures test RMSE minus train RMSE; larger gaps indicate overfitting.}
\label{tab:generalization}
\begin{tabular}{lccc}
\toprule
\textbf{Model} & \textbf{Train RMSE} & \textbf{Test RMSE} & \textbf{Gap} \\
\midrule
NXRO-Res & 0.474°C & 0.567°C & 0.093°C \\
NXRO-Graph & 0.536°C & 0.583°C & 0.047°C \\
NXRO-Attentive & 0.542°C & 0.586°C & 0.044°C \\
NXRO-Linear & 0.543°C & 0.588°C & 0.045°C \\
NXRO-NeuralODE & 0.522°C & 0.588°C & 0.066°C \\
XRO & 0.493°C & 0.605°C & 0.112°C \\
NXRO-RO+Diag & 0.490°C & 0.618°C & 0.128°C \\
NXRO-Transformer & 0.678°C & 0.651°C & $-$0.027°C \\
\bottomrule
\end{tabular}
\end{table}

Several patterns emerge from this analysis. Models with learned sparsity or restricted interactions (Graph, Attentive, Linear) exhibit the smallest gaps, suggesting that architectural constraints provide effective regularization. The physics-structured models (XRO, NXRO-RO+Diag) exhibit larger gaps despite achieving excellent training performance, indicating that the polynomial basis may capture training-period patterns that do not persist.

The residual model occupies an interesting middle ground: it achieves the best training performance among the neural variants but also exhibits a larger gap than Graph or Attentive. This suggests that the MLP residual has sufficient capacity to overfit somewhat, but not catastrophically. The balance between expressiveness and regularization appears well-calibrated for this problem.

The transformer's negative gap (test better than train) reflects its failure to learn meaningful dynamics from either dataset. When a model cannot fit the training data, the train-test gap becomes uninformative about generalization.

\subsection{Transfer Learning Analysis}

We investigate whether pre-training on synthetic climate data can improve generalization~\cite{pan2010transfer}. The two-stage training procedure first trains on CMIP6-derived~\cite{eyring2016cmip6} climate indices spanning thousands of simulation years, then fine-tunes on the observational training period. Figure~\ref{fig:two_stage} compares single-stage (observations only) and two-stage (pre-training plus fine-tuning) training across model variants.

\begin{figure}[htbp]
    \centering
    \includegraphics[width=0.48\textwidth]{comparison_single_vs_two_stage_rmse.png}
    \hfill
    \includegraphics[width=0.48\textwidth]{comparison_single_vs_two_stage_acc.png}
    \caption{Comparison of single-stage versus two-stage training. Left: Test RMSE, where bars below zero indicate two-stage improvement. Right: Test ACC, where bars above zero indicate two-stage improvement.}
    \label{fig:two_stage}
\end{figure}

The effect of pre-training varies dramatically across architectures. For the physics-structured model (NXRO-RO+Diag), two-stage training proves essential: single-stage training produces unstable optimization with divergent forecasts, while two-stage training yields competitive performance (RMSE $\sim$0.63°C). The pre-trained weights provide an initialization in a stable region of parameter space that gradient descent can refine.

For the pure neural ODE, pre-training provides modest but consistent benefits. The model learns general features of climate variability from the abundant synthetic data, providing a better starting point than random initialization. The improvement is smaller than for the physics model because the MLP has fewer parameters and is easier to optimize from scratch.

For the residual model (NXRO-Res), pre-training actually degrades performance, increasing test RMSE from 0.568°C to 0.619°C. This counterintuitive result arises because the synthetic climate data, while realistic, contains systematic biases relative to observations. The MLP learns to correct for these biases during pre-training, but these corrections are inappropriate for the real system. During fine-tuning, the model must unlearn these biased corrections, a process that may leave residual artifacts.

This asymmetry suggests a principle for transfer learning in hybrid models: pre-training benefits components that encode general structure (like the physics model's polynomial coefficients) more than components that capture dataset-specific corrections (like the MLP residual). When the pre-training distribution differs from the target distribution, flexible components may learn spurious patterns that are difficult to unlearn.

\subsection{Selective Freezing Experiments}

The freezing experiments described in Section~\ref{sec:neural_ode} reveal which XRO components are robust enough to fix at their closed-form values versus which benefit from gradient-based refinement.

Freezing the RO coupling coefficients while training the linear operator (NXRO-FixRO) yields test RMSE of 0.602°C, slightly better than full XRO (0.605°C). This suggests that XRO's core physics---the polynomial coupling between temperature and thermocline depth---captures robust dynamics that do not require refinement, while the linear operator benefits from end-to-end optimization that accounts for forecast error accumulation.

Freezing both RO and diagonal terms while training only the linear operator (NXRO-FixNL) achieves the best training RMSE (0.470°C) among all variants but degrades to 0.602°C on the test set. The frozen nonlinearities constrain the optimization to adjust only linear coupling strengths, preventing the overfitting that afflicts the fully-trainable physics model.

Adding an MLP residual to the frozen physics (NXRO-ResidualMix-FixRO) yields test RMSE of 0.598°C, providing no improvement over the simpler residual model. This suggests that XRO's polynomial nonlinearities and the MLP residual capture overlapping aspects of the dynamics; combining both provides redundancy rather than complementary information.

\subsection{Comparison Across Evaluation Regimes}

The rankings change substantially when evaluated on training data versus held-out test data, as shown in Table~\ref{tab:ranking_changes}. Models that rank highly in-sample often perform worse out-of-sample, and vice versa.

\begin{table}[htbp]
\centering
\caption{Ranking comparison between in-sample and out-of-sample evaluation. Dashes indicate the model was not among the top performers in that evaluation regime.}
\label{tab:ranking_changes}
\begin{tabular}{lcc}
\toprule
\textbf{Model} & \textbf{In-Sample Rank} & \textbf{Out-of-Sample Rank} \\
\midrule
NXRO-ResidualMix & 1 & --- \\
NXRO-RO+Diag & 2 & 4 \\
XRO & 3 & 7 \\
NXRO-Res & --- & 1 \\
NXRO-Graph & --- & 2 \\
NXRO-Attentive & --- & 3 \\
\bottomrule
\end{tabular}
\end{table}

The ResidualMix variant, which combines XRO's full physics structure with an MLP residual, achieves the best in-sample performance but fails to appear in the top out-of-sample rankings. Its high capacity allows it to fit training trajectories precisely, but this flexibility becomes a liability when the test period exhibits different dynamics.

Conversely, the simpler hybrid models (Res, Graph, Attentive) that rank highest out-of-sample achieve only moderate in-sample performance. Their constrained architectures prevent perfect training fit but yield representations that generalize to unseen data.

This pattern underscores the importance of out-of-sample evaluation for model selection in climate forecasting. In-sample metrics, while easier to compute, systematically favor complex models that overfit. Rigorous temporal holdout is essential for identifying models with genuine predictive skill.

\subsection{Forecast Skill Across Lead Times}

Figure~\ref{fig:all_variants} displays the complete comparison across all model variants, showing both ACC and RMSE for training and test periods. Several lead-time-dependent patterns emerge.

\begin{figure}[htbp]
    \centering
    \includegraphics[width=0.95\textwidth]{all_variants_comparison_combined_out_of_sample.png}
    \caption{Performance comparison across all model variants. Top panel: ACC for train (blue) and test (orange) periods. Bottom panel: RMSE for train and test periods. The gap between train and test performance indicates overfitting tendency.}
    \label{fig:all_variants}
\end{figure}

At short lead times (1--3 months), all models achieve high skill (ACC $>$ 0.85, RMSE $<$ 0.4°C), and the differences between variants are small. ENSO persistence provides strong predictability at these horizons, and even simple linear models capture most of the signal.

At medium lead times (4--9 months), model differences become pronounced. The residual model maintains higher ACC than XRO, with the advantage growing from leads 4--6 (within the spring barrier) where traditional models struggle. The learned residual appears to capture dynamics that help maintain skill through this challenging period.

At long lead times (12--21 months), the residual and graph models maintain moderate skill (ACC $\sim$ 0.4--0.5) while XRO and the physics-structured variants degrade more rapidly (ACC $\sim$ 0.3--0.4). The polynomial nonlinearities that help XRO fit training trajectories may encode amplitude-dependent patterns specific to the 1979--2001 period that do not persist into 2002--2022.

\subsection{Summary and Recommendations}

Our experimental evaluation yields several actionable findings for practitioners.

The residual neural model (NXRO-Res) represents the recommended default for operational ENSO forecasting. Its combination of physics-informed linear structure and flexible neural residual achieves the best out-of-sample skill while remaining computationally efficient and straightforward to train. The architecture is robust to initialization and hyperparameter choices, requiring minimal tuning beyond the defaults described in Section~\ref{sec:neural_ode}.

For applications requiring interpretable coupling patterns, the graph-based model (NXRO-Graph) offers a compelling alternative. Its sparse connectivity structure can be inspected to understand which climate modes interact most strongly, and its performance trails the residual model by only 0.016°C in RMSE.

The transformer architecture should be avoided for climate forecasting with current data volumes. Its failure despite architectural sophistication illustrates that modern deep learning techniques do not automatically transfer to small-data scientific domains. Physics-informed inductive biases are not merely helpful but essential when training data is limited.

Transfer learning from climate model simulations provides benefits for some architectures (physics-structured models, pure neural ODEs) but can harm others (residual models). This finding aligns with recent work on combined dynamical-deep learning ENSO forecasts~\cite{wang2025combined}, which demonstrates that the optimal integration of physics and machine learning depends critically on the model architecture. Practitioners should evaluate two-stage training on held-out data before deployment rather than assuming pre-training helps.

Finally, out-of-sample evaluation is essential for model selection. In-sample metrics systematically favor complex models that overfit, providing misleading guidance about operational performance. Any evaluation of ENSO forecast systems should employ rigorous temporal holdout spanning multiple ENSO cycles.
